\cotefile{93}
\scanpage{151}
donnée d'un hom de faisceaux
\[
 P \xrightarrow{\lambda_{\deleted{\unsure{P|E}}}} P^\infty_{S/T}(E)/\delta E \tag*{(26)}
\]
satisfaisant la condition
\[
 \deleted{\lambda(x+\unsure{\xi|\zeta})=\lambda(x)+d^\infty_{S/T}(\xi)\qquad x\in P(\cdot),\ \xi\in E(\cdot).} \tag*{\deleted{(27)}}
\]
\[
 \lambda(x+\xi)=\lambda(x)+p(d^\infty_{S/T}(\xi))\qquad x\in P(\cdot),\ \xi\in E(\cdot). \tag*{(27)}
\]
Il faut encore exprimer que la trivialisation du $P^\infty_{S/T}(E)/E$-torseur obtenu est compatible avec les stratifications. Comme $P$ est \deleted{le} obtenu des sections horizontales de $P^\infty_{S/T}(P)$, on voit que \deleted{\unsure{cette|une}}
% slanted
\margin{plutôt lemme à dégager pour un hom de $P^\infty_{S/T}(P)$ dans un torseur stratifié quelc.}
\deleted{condition \unsure{de|du}} \inserted{cela} signifie qu'on a
\[
 \lambda(P)\subset\bigl(P^\infty_{S/T}(E)/E\bigr)^{\natural} \,, \tag*{(28)}
\]
où le \deleted{\unsure{donné|demi}} signe $\natural$ désigne le sous faisceau des sections horizontales.

\deleted{En} \uline{En résumé}, \deleted{\unsure{une struc|une donnée}} si $P$ est un $E$-torseur, une stratification sur $E$ rel. à $T$, compatible avec la stratification $\delta=\delta^\infty_E$ donnée sur $E$, équivaut à un hom de faisceaux (26), satisfaisant (27) et (28) ; cette dernière condition pouvant d'ailleurs s'exprimer en précisant la donnée (26) par
% Two vertical emphasis strokes alongside the four preceding lines.
\[
 P\xrightarrow{\lambda}\bigl(P^\infty_{S/T}(E)/\delta E\bigr)^\natural . \tag*{(26$'$)}
\]
Cette description est fonctorielle en $E$, et compatible aux changements de base.
% archive sheet number: 150
